\section{Slide y estrategia de evidencia visual}

\subsection{Archivo estructurado slide-first}
El expositor, al narrar el proceso de optimización de trace, agent y DSPy, se apoya siempre en imágenes de slide estructuradas para resaltar el pipeline completo. En el proceso de edición seguimos el principio slide-first: primero capturamos los slides estáticos (como las tres categorías de estrategias, la Constitución APAC, la arquitectura de DSPy y la latent optimization), y después, donde sea necesario, agregamos frames dinámicos, usando evidencia visual para anclar cada tramo de la lógica didáctica. La siguiente tabla enumera el material visual principal que se usa hasta ahora y su valor didáctico.
\begin{table}[H]
\centering
\begin{tabular}{lp{8cm}}
\toprule
Slide/Frame & Contenido principal y uso \\
\midrule
frame-solution-categories.jpg & El slide de «capas» donde aparecen las tres estrategias (Scaling/Hybrid/Compound AI), que ayuda al lector a entender el trade-off de costo/control entre los distintos métodos (00:09:10--00:09:18). \\
frame-apac-constitution.jpg & La tabla de las cuatro dimensiones de la APAC Constitution, que explica cómo se vinculan las métricas discriminativas con la estrategia de preguntas del agent (00:16:08--00:16:22). \\
frame-dspy-architecture.jpg & El diagrama de arquitectura de DSPy, que muestra el ciclo cerrado entre latent cost, solver y plan (00:32:25--00:32:45). \\
frame-latent-opt.jpg & El diagrama diferenciable que describe el pipeline de latent optimization, usado para explicar la iteración de dos niveles de Green Descent (00:45:05--00:45:25). \\
\bottomrule
\end{tabular}
\caption{Slides/frames usados en esta sección y su valor didáctico respectivo}
\end{table}
\begin{knowledgebox}{Sinergia entre slide y frame}
Los slides se usan para presentar fórmulas estructuradas, hojas de ruta y tablas de métricas; los frames se usan para capturar el estado del expositor mientras explica, el movimiento de un diagrama esquemático o un ejemplo numerical llamativo. Mientras haya un slide disponible, se prioriza el slide, y el frame sirve como complemento dinámico o para llenar puntos ciegos de la explicación; juntos sostienen una narrativa didáctica continua.
\end{knowledgebox}

\subsection{Notas de tiempo de los frames clave}
Cada ilustración registra estrictamente su intervalo de tiempo correspondiente, colocado en el pie de figura o en el footer, para que el lector pueda usar el subtitle y volver a los detalles. Por ejemplo, trace curve, la estructura APAC y el latent optimizer se marcan, cada uno en su propia subsección, con el formato \texttt{00:xx:xx--00:yy:yy}. Si en el futuro se agregan nuevos slides, es indispensable primero comparar con los subtitles para encontrar el frame más completo que se muestre, y anotar con claridad el intervalo de tiempo.
\begin{warningbox}{La marca de tiempo de la imagen no se puede omitir}
La falta de una marca de tiempo impide que el lector confirme el origen de la visual evidence, sobre todo cuando varias imágenes contiguas tratan temas distintos. Cada figure debe dar el intervalo de subtítulos en la misma página, y de ser necesario, repetirlo en el caption.
\end{warningbox}

\subsection{Resumen de la sección}
La forma de organización slide-first le da al pipeline abstracto anclas visuales, y junto con los ejemplos dinámicos del frame y las marcas de tiempo precisas, permite que el lector, al leer el texto, regrese rápidamente a la escena original de la explicación en el video.

